% -----------------------------------------------------------------------------------------------------
\section{引言：范围、重叠规则与阅读深度}\label{sec:front-intro}

本书第二至第六部分的证明建立在少数几条基础性陈述之上：钳夹面在哪里，以及如何用 Landau 方程与 Coleman--Norton 图像找到它们（\cref{sec:ir-landau}）；OPE 在什么意义下成立（\cref{thm:dis-ope}、\cref{thm:power-ope-convergence}）；区域方法作为计算工具与 SCET 模式表的地位（\cref{sec:ir-regions}）；Glauber 区为什么在单举 Drell--Yan 中相消而在别处未必相消（\cref{thm:dy-factorization}，\cref{sec:break-mpi}）；非全局可观测量中的对数结构（\cref{sec:jets-ngl}）；以及幂次修正如何组织（\cref{prop:power-nlp-endpoint}）。本章要回答的问题是：自 2010 年前后以来，这些基础中哪些发生了变化，朝哪个方向变化。答案大致分为三类。第一类是把原先只有启发式论证的结论变为有完整书面证明的定理，最突出的是 Collins 的物理区钳夹判据（\cref{thm:front-collins-pinch}；由于证明只见于预印本，本书仍记为 [A]）。第二类是把被普遍使用的一般陈述证伪：“区域 $=$ Newton 多面体的下表面”在 Minkowski 运动学中不成立（\cref{prop:front-hidden-regions}），$\MSbar$ 部分子分布并非在所有标度上非负（\cref{prop:front-msbar-negative}），1-jettiness 的标准指数化形式在强子对撞中不成立（\cref{prop:front-jettiness-cvl}），简单 Regge 极点因子化在次次领头对数精度失效（\cref{prop:front-regge-pole}）。第三类是在原有框架之内的精细化：超领头对数的全阶重求和、次领头幂次的端点因子化、超出领头幂次的 TMD 算符展开、无重复计数的双部分子散射方案、由结构函数确定光子分布等。

\paragraph{重叠规则。}前面各章已经陈述过的核心结论在这里从不重述。本章只给出前沿的细化，以 \texttt{front} 为前缀的标签表示，并在首次出现时注明所引核心陈述的标签与等级。\cref{tab:front-map}列出这种对应。

\paragraph{阅读深度、预印本与分歧。}本章引用的文献很多，我们对它们的阅读深度并不相同。对读过全文的文献（例如 Collins 关于 Landau 条件的论文\cite{Collins:2020euz}，Gardi、Herzog、Jones、Ma 与 Schlenk 关于在壳展开的论文\cite{Gardi:2022khw}），我们给出定理编号与公式位置；只读过摘要的文献，相应的转述一律写成“据其摘要”或标注“（据摘要）”，并且不据此作出超出摘要的断言。截至 2026 年 10 月仍只有 arXiv 版本的文献在首次引用时标注“（预印本）”，在依赖它们的「地位与缺口」段中再次说明；本章涉及的有\cite{Collins:2020euz}（预印本）、\cite{Frob:2016mzv}（预印本）、\cite{Guida:2011np}（预印本）、\cite{Beneke:2026miy}（预印本）、\cite{Gamberg:2022lju}（预印本）、\cite{Bodwin:2025zzi}（预印本）、\cite{Anastasiou:2025cvy,Anastasiou:2026kpm,Anastasiou:2026jta}（预印本）与\cite{Becher:2026zon,Becher:2026iim}（预印本）。不同研究组之间的分歧，我们并列给出各自结论的确切适用范围，不作裁决。

\paragraph{全局性保留。}本章所有“严格”的结果都是微扰论（在每一固定圈数）或模型中的陈述；唯一的非微扰正面结论是共形场论中的光线 OPE（\cref{prop:front-lightray-ope}(b)），其余的非微扰陈述都是反例。本章没有任何一条结论控制 QCD 中因子化的非微扰有效性；后者仍由第一章的全局假设（例如\cref{asm:intro-bridge}）承担。

\paragraph{约定。}除另有说明外，本章沿用全书约定（光锥分量 $v^\pm=(v^0\pm v^3)/\sqrt2$，协变导数 $D_\mu=\partial_\mu-\ii gA_\mu^at^a$）。三处例外在用到时说明：\cref{sec:front-landau}沿用 Collins 的积分记号（分母 $A_j$、梯度 $D_j$、二阶导数矩阵 $E_j$），它与\cref{thm:ir-landau-momentum}的记号不同；\cref{sec:front-glauber}与\cref{sec:front-nlp}所引的 SCET 算符采用\cref{def:scet-modes}的约定（$n\cdot\bar n=2$）；喷注文献中的多 Wilson 线软函数 $S_i=\mathbf P\exp(\ii g_s\int_0^\infty\dd s\,n_i\cdot A^a_s\,\mathbf T^a_i)$\cite{Becher:2016mmh}所带的耦合常数符号与本书相反（换算到本书约定：$g\to-g$），由于本章只引用由它们得到的、对 $g$ 为偶的结果，这一差别不影响任何公式。

\begin{table}[htbp]
\centering
\caption{前沿主题、核心陈述与本章的细化。中间一列是前面各章（或附录）中的核心陈述及其等级；右列是本章对应的细化条目。}\label{tab:front-map}
\small
\renewcommand{\arraystretch}{1.12}
\begin{tabularx}{\textwidth}{@{}p{2.5cm}>{\raggedright\arraybackslash}X>{\raggedright\arraybackslash}X@{}}
\toprule
主题（节） & 核心陈述（等级） & 本章的细化 \\
\midrule
Landau 判据、Coleman--Norton（\cref{sec:front-landau}） & \cref{thm:ir-landau-momentum}（[R, $R_{\mathrm{PT}}$]）、\cref{rem:ir-landau-status}（[A]）、\cref{prop:ir-coleman-norton}（[R, $R_{\mathrm{PT}}$]）；\cref{lem:appD-gordan}（数学） & \cref{def:front-trap,lem:front-min-norm,thm:front-collins-pinch,cor:front-coleman-norton,prop:front-anomalous,prop:front-landau-not-singular,prop:front-ir-forest,prop:front-coordinate-space,prop:front-local-ir} \\
OPE（\cref{sec:front-ope}） & \cref{thm:dis-ope}（[A]）、\cref{thm:power-ope-convergence}（[R, $R_{\mathrm{PT}}$]） & \cref{prop:front-ope-scalar,prop:front-ope-ym,prop:front-flow-renorm,conj:front-ope-convergence} \\
区域方法（\cref{sec:front-regions}） & \cref{def:ir-regions}、\cref{thm:ir-regions-euclid}（[R, $R_{\mathrm{PT}}$]）、\cref{prop:ir-regions-minkowski}（[P]）；\cref{def:scet-zero-bin} & \cref{lem:front-jantzen,def:front-facets,prop:front-lp-facets,prop:front-onshell-facets,prop:front-commuting,prop:front-hidden-regions,conj:front-mor-minkowski} \\
Glauber（\cref{sec:front-glauber}） & \cref{prop:scet-glauber-operators}（[A]）、\cref{thm:dy-factorization}（[A]）、\cref{prop:break-spacelike,prop:break-fss,prop:break-jet-veto}（[V]）、\cref{prop:break-mass-singularity}（[C]） & \cref{prop:front-glauber-bfkl,prop:front-syz,conj:front-mpi-violation,prop:front-heavy-fsi} \\
超领头对数（\cref{sec:front-sll}） & \cref{prop:jets-sll}（[P]）、\cref{thm:jets-ngl-factorization}（[A]） & \cref{prop:front-sll-all-orders,prop:front-sll-closed-form,prop:front-restoration,prop:front-jettiness-cvl} \\
次领头幂次（\cref{sec:front-nlp}） & \cref{prop:thr-nlp}（[P]）、\cref{prop:power-nlp-endpoint}（[P]） & \cref{thm:front-hgg,prop:front-ggh,prop:front-nlp-threshold,prop:front-nlp-ll,prop:front-nlp-dis,conj:front-nlp-beyond-ll} \\
次领头幂次 TMD（\cref{sec:front-nlp-tmd}） & \cref{thm:tmd-factorization}（[A]） & \cref{conj:front-glauber-nlp,prop:front-tmd-nlp-operator,prop:front-sidis-nlp,prop:front-nlp-tmd-oneloop} \\
正定性（\cref{sec:front-positivity}） & \cref{conj:pdf-positivity}（[C]） & \cref{prop:front-msbar-negative,prop:front-positivity-large-scale} \\
BRST 分类（\cref{sec:front-brst}） & \cref{prop:appJ-local-operators}（[R, $R_{\mathrm{PT}}$]）、\cref{rem:appJ-nonlocal}、\cref{prop:appL-renormalized}（[R, $R_{\mathrm{PT}}$]） & \cref{prop:front-brst-bilocal,conj:front-brst-multilocal} \\
小 $x$（\cref{sec:front-smallx}） & \cref{thm:smallx-bfkl,thm:smallx-bk}（[A]） & \cref{prop:front-nlo-bk,prop:front-bk-collinear,prop:front-regge,prop:front-regge-pole,prop:front-smallx-tmd,prop:front-smallx-helicity} \\
双部分子散射（\cref{sec:front-dps}） & \cref{def:break-dpdf}、\cref{thm:break-double-dy}（[A]）、\cref{conj:break-dps-coloured}（[C]） & \cref{prop:front-dps-scheme,prop:front-rapidity-renorm,prop:front-dpd-evolution} \\
喷注与非全局对数（\cref{sec:front-jets}） & \cref{thm:jets-ngl-factorization,thm:jets-bms}（[A]）、\cref{prop:jets-ngl-beyond}（[P]） & \cref{prop:front-ngl,prop:front-jet-substructure,prop:front-lightray-ope} \\
重夸克偶素（\cref{sec:front-quarkonium}） & \cref{conj:heavy-nrqcd-production}（[C]）、\cref{prop:heavy-quarkonium-fragmentation}（[A]） & \cref{prop:front-nrqcd-ir} \\
QED 与电弱（\cref{sec:front-qed}） & \cref{prop:ir-bn-failure}（[V]）、\cref{prop:coll-luxqed}（[A]） & \cref{thm:front-luxqed,prop:front-lepton-pdf,prop:front-ew-pdf,prop:front-ew-violation} \\
\bottomrule
\end{tabularx}
\end{table}

